% paper/sections/08d_gfm.tex
%
% 付録 D.4  Hermite 場延長法（HFE）
%   本文§8.4から付録Dへ移動 (2026-04-05)
%
%   Aslam (2004) の Extension PDE の定常解を CCD Hermite データから
%   直接構成し，界面越し場延長を O(h^6) で実現する．
% ═══════════════════════════════════════════════════════════════════════════════
\section{\texorpdfstring{Hermite 場延長法（HFE）：$\Ord{h^6}$ 界面越し場延長}{Hermite 場延長法（HFE）：O(h⁶) 界面越し場延長}}
\label{sec:field_extension}
\label{sec:hfe}% canonical label
\label{sec:extension_pde}% backward-compat: §1–§9 cross-refs
\label{sec:gfm}% backward-compat: §10-§12 cross-refs — 削除条件: sections/{10,11,12,13}*.tex で \\ref{sec:gfm} が Grep ヒット0件になった後
% ═══════════════════════════════════════════════════════════════════════════════

本節では，\textbf{Hermite 場延長法}（Hermite Field Extension, 以下 \textbf{HFE}）を定式化する．
HFE は Aslam (2004)~\cite{Aslam2004} の Extension PDE の定常解を，
CCD の Hermite データ $(f, f', f'')$ から仮想時間反復を経由せず直接構成し，
$\Ord{h^6}$ 精度の界面越し場延長を実現する手法である．

\noindent\textbf{ソルバにおける位置づけ．}
低密度比（$\rho_l/\rho_g \leq 5$）の smoothed Heaviside 一括解法では，
圧力場 $p$ は $\varepsilon$-regularized な滑らかな場であり，
HFE は不要である（むしろ有害；§~\ref{sec:ppe_discretization_choice} で実証）．
高密度比（$\rho_l/\rho_g > 5$）では\textbf{分相 PPE 解法}
（§~\ref{sec:split_ppe_motivation}，§~\ref{sec:split_ppe_recovery} で実証）を採用し，
各相で独立に PPE を解く．圧力場が Young--Laplace ジャンプ $[p]_\Gamma = \sigma\kappa$ を
陽に持つため，HFE による界面越し場延長が\textbf{必須}となる．
本節ではこの分相 PPE 解法の\textbf{基盤技術}として HFE を定式化する．

\paragraph{なぜ分相 PPE 解法に界面越し場延長が必要か.}
分相 PPE 解法では各相で定数密度の PPE を独立に解くため，
圧力場が Young--Laplace 則 $[p]_\Gamma = \sigma\kappa$ に従う
不連続を持つ（式~\eqref{eq:young_laplace_jump}）．
CCD 法は同一ステンシル $(i-2,\ldots,i+2)$ で $\Ord{h^6}$ の微分を求解するが，
界面を跨ぐステンシルは Gibbs 振動を引き起こし精度が著しく低下する．

HFE はこの問題を根本的に解消する：
ソース相（例：液相）の場の値を界面法線方向にターゲット相へ $\Ord{h^6}$ 精度で延長し，
CCD が界面近傍でも滑らかな場に対して演算できるようにする．
具体的な適用箇所は以下の通りである：
\begin{enumerate}[noitemsep]
  \item \textbf{圧力場の前処理：}
        PPE 求解前に $p^n$ を界面越しに延長し，IPC Predictor の $-\bnabla p^n$ 項から
        Gibbs 振動を除去する．
  \item \textbf{圧力増分の後処理：}
        PPE で得られた $\delta p$ を界面越しに延長し，速度補正
        $\bu^{n+1} = \bu^* - (\Delta t/\rho)\,\bnabla(\delta p_{\mathrm{ext}})$ の精度を維持する．
  \item \textbf{変密度 PPE の精度改善：}
        界面近傍で $1/\rho$ が急変する領域において，延長済みの滑らかな場を用いることで
        CCD の離散化誤差を低減する（変密度 PPE の収束性限界については
        §~\ref{sec:varrho_ppe_limitation} で定量的に記録する）．
\end{enumerate}

\subsubsection{背景：Extension PDE の定常解と HFE の着想}

スカラー場 $q$（圧力増分 $\delta p$ 等）を液相側（$\phi < 0$）から気相側（$\phi \ge 0$）へ滑らかに延長する
偏微分方程式は：
\begin{equation}
  \frac{\partial q}{\partial \tau} + \mathrm{sgn}(\phi)\,\hat{\bm{n}}\cdot\bnabla q = 0
  \label{eq:extension_pde_main}
\end{equation}
\phantomsection\label{eq:extension_pde}% alias used by §3b cross-refs
ここで $\tau$ は仮想時間（次元は長さ $[L]$；特性速度 $|\hat{\bm{n}}| = 1$ により $\Delta\tau \sim h$），
$\hat{\bm{n}} = \bnabla\phi/|\bnabla\phi|$ は界面法線ベクトルであり，
$\mathrm{sgn}(\phi) = +1$（気相）$/-1$（液相）が延長方向を定める．
液相側（ソース相）の値を固定したまま，気相側（ターゲット相）を仮想時間前進させることで
$\bnabla q \cdot \hat{\bm{n}} = 0$（法線方向勾配ゼロ）の定常解へ収束する．

\subsubsection{CCD との整合性}

最近接点 Hermite 補間（§~\ref{sec:extension_pde_disc}）により構成された延長後の場 $q_{\mathrm{ext}}$ は，
定常解 $\bnabla q \cdot \hat{\bm{n}} = 0$ を $\Ord{h^6}$ 精度で満たすため，
界面近傍で $C^\infty$ に近い滑らかな分布を持つ（図~\ref{fig:extension_schematic}）．
この滑らかな場 $q_{\mathrm{ext}}$ に対して CCD 演算子 $D^{(1)}, D^{(2)}$ を適用することで：
\begin{itemize}
  \item 界面直近まで $\Ord{h^6}$ 精度の圧力勾配 $\bnabla(\delta p)$ が得られる
  \item 速度補正 $\bu^{n+1} = \bu^* - (\Delta t/\rho)\,\bnabla(\delta p)$ の離散化誤差が低減される
  \item Balanced-Force 条件（§~\ref{sec:balanced_force}）と整合した圧力勾配評価が実現される
\end{itemize}

\begin{figure}[htbp]
\centering
\begin{tikzpicture}[scale=0.85]
  % 延長前
  \begin{scope}[xshift=-50mm]
    \node[font=\small\bfseries, text=navy] at (0,2.8) {延長前：$\delta p$ に不連続};
    \draw[->] (-2.5,0) -- (2.5,0) node[right,font=\small]{$x$};
    \draw[->] (0,-0.5) -- (0,2.5) node[above,font=\small]{$\delta p$};
    \draw[very thick, blue2, domain=-2.2:0] plot (\x, {1.8});
    \draw[very thick, red2, domain=0:2.2] plot (\x, {0.3});
    \draw[dashed, gray] (0,-0.3) -- (0,2.3);
    \node[font=\scriptsize] at (0,-0.5) {$\Gamma$};
    \node[font=\scriptsize, blue2] at (-1.2, 2.1) {液体};
    \node[font=\scriptsize, red2] at (1.2, 0.6) {気体};
    \draw[thick, orange2, decorate, decoration={zigzag, amplitude=3pt, segment length=5pt}]
      (-0.3, 1.0) -- (0.3, 1.0);
    \node[font=\scriptsize, orange2] at (0, 0.6) {Gibbs};
  \end{scope}
  % 矢印
  \draw[-Stealth, very thick, navy] (-1.5, 1.2) -- (1.0, 1.2)
    node[midway, above, font=\small, navy]{Hermite 延長};
  % 延長後
  \begin{scope}[xshift=50mm]
    \node[font=\small\bfseries, text=navy] at (0,2.8) {延長後：$\delta p_\mathrm{ext}$ は滑らか};
    \draw[->] (-2.5,0) -- (2.5,0) node[right,font=\small]{$x$};
    \draw[->] (0,-0.5) -- (0,2.5) node[above,font=\small]{$\delta p_\mathrm{ext}$};
    \draw[very thick, blue2, domain=-2.2:2.2] plot (\x, {1.8});
    \draw[dashed, gray] (0,-0.3) -- (0,2.3);
    \node[font=\scriptsize] at (0,-0.5) {$\Gamma$};
    \node[font=\scriptsize, blue2] at (1.2, 2.1) {延長済み};
    \node[font=\scriptsize, teal2] at (0, 0.4) {CCD $\bnabla(\delta p_\mathrm{ext})$: $\Ord{h^6}$};
  \end{scope}
\end{tikzpicture}
\caption{最近接点 Hermite 補間による圧力場の界面越し延長（概念図）．
左：CSF 表面張力による Laplace ジャンプが CCD ステンシルを跨ぐと Gibbs 振動が発生．
右：最近接点 Hermite 補間で液体側の値を気体側へ $\Ord{h^6}$ 精度で延長した後は，
CCD $\bnabla(\delta p)$ が $\Ord{h^6}$ 精度を界面直近まで維持する．}
\label{fig:extension_schematic}
\end{figure}

\subsubsection{HFE の離散化：最近接点 Hermite 補間}
\label{sec:extension_pde_disc}

式~\eqref{eq:extension_pde_main} の定常解は最近接点延長
$q_\mathrm{ext}(\bm{x}) = q\bigl(\bm{x}_\Gamma(\bm{x})\bigr)$ に他ならない
（Aslam~(2004)~\cite{Aslam2004}, §2；$\bnabla q \cdot \hat{\bm{n}} = 0$ の
特性曲線解が最近接点の値の伝播と一致する）．
ここで $\bm{x}_\Gamma$ は $\bm{x}$ から界面 $\Gamma$ への最近接点である．
本稿ではこの定常解を，Extension PDE の仮想時間前進を経由せず，
CCD の Hermite データ $(f, f', f'')$ を用いた直接補間により構成する．

\paragraph{最近接点の計算.}
ターゲット側格子点 $\bm{x}$（$\phi(\bm{x}) \geq 0$）に対し，
法線 $\hat{\bm{n}} = \bnabla\phi/|\bnabla\phi|$ を CCD $D^{(1)}$（$\Ord{h^6}$）で評価し，
最近接点を
\begin{equation}
  \bm{x}_\Gamma = \bm{x} - \phi(\bm{x})\,\hat{\bm{n}}(\bm{x})
  \label{eq:closest_point}
\end{equation}
で求める．$\phi$ が再初期化により $|\bnabla\phi| \approx 1$（SDF 条件）を満たすとき，
$\bm{x}_\Gamma$ の位置精度は CCD の $\bnabla\phi$ 精度で律速され $\Ord{h^6}$ となる．

\paragraph{CCD Hermite 5次補間.}
CCD は各格子点 $x_i$ で $(f_i, f'_i, f''_i)$ を同時に返す．
隣接2点 $x_a, x_b$（$h = x_b - x_a$，いずれもソース側）の
6個のデータから5次 Hermite 多項式を一意に構成する：
\begin{equation}
  P(\xi) = \sum_{k=0}^{5} c_k \,\xi^k, \qquad \xi = \frac{x - x_a}{h} \in [0,1]
  \label{eq:hermite5}
\end{equation}
拘束条件は $P^{(j)}(0) = h^j f_a^{(j)}$，$P^{(j)}(1) = h^j f_b^{(j)}$（$j = 0,1,2$）の6個であり%
\footnote{$h^j$ は $\xi = (x-x_a)/h$ への無次元化に伴うスケール因子：
$P'(\xi) = h f'(x)$，$P''(\xi) = h^2 f''(x)$．}，
係数 $c_0,\ldots,c_5$ は閉形式で得られる．
補間誤差は
\begin{equation}
  |P(x) - f(x)| \leq \frac{\|f^{(6)}\|_\infty}{6!}\,|(x-x_a)^3(x-x_b)^3| = \Ord{h^6}
  \label{eq:hermite5_error}
\end{equation}
であり，CCD 微分精度と同じ $\Ord{h^6}$ を保つ．

\paragraph{片側外挿.}
$\bm{x}_\Gamma$ がソース側最端セル $[x_i, x_{i+1}]$ の外側
（$x_{i+1}$ がターゲット側）に位置する場合，
ソース側のみの2点 $[x_{i-1}, x_i]$ を用いて外挿する．
外挿距離は高々 $h$ であり，式~\eqref{eq:hermite5_error} から
$|(2h)^3 \cdot h^3| = 8h^6$ となるため $\Ord{h^6}$ が維持される．

\paragraph{2次元への拡張.}
テンソル積逐次補間を用いる．
$\bm{x}_\Gamma = (\xi, \eta)$ に対し，$y$ 方向各行で $x$ 方向 Hermite 補間を行い，
得られた値・微分を $y$ 方向に再度 Hermite 補間する．
必要な CCD 演算は $q_x, q_{xx}$（$x$ 軸），$q_y, q_{yy}$（$y$ 軸），
$q_{xy}, q_{xyy}$（混合微分）の計4回であり，
仮想時間反復（$n_\mathrm{ext} \times 2$ 軸 $= 10$ 回）より少ない．

\begin{tcolorbox}[defbox, title={補足：仮想時間前進による離散化との比較}]
Aslam~(2004)~\cite{Aslam2004} の元論文では，式~\eqref{eq:extension_pde_main} を
1次風上差分＋前進 Euler で仮想時間反復する手法が提案されている．
この手法は実装が単純だが，風上差分の数値拡散により延長帯の精度が $\Ord{h}$ に制限される．
HFE は同一の定常解を CCD データから直接構成することで
精度を $\Ord{h^6}$ に引き上げると同時に仮想時間反復を排除する．
\end{tcolorbox}

\subsubsection{HFE のアルゴリズムへの組み込み}

時刻ループ内での適用箇所（§~\ref{sec:algorithm}）：
\begin{enumerate}
  \item \textbf{PPE 求解前（Step 5c）：} 前時刻圧力 $p^n$ に HFE を適用して $p^n_{\mathrm{ext}}$ を構築し，
        IPC Predictor の $-\bnabla p^n$ 項に使用する
  \item \textbf{Corrector 前（Step 7 前）：} 圧力増分 $\delta p$ に HFE を適用して $\delta p_{\mathrm{ext}}$ を構築し，
        速度補正 $\bu^{n+1} = \bu^* - (\Delta t/\rho)\,\bnabla(\delta p_{\mathrm{ext}})$ に使用する
\end{enumerate}

\medskip
\noindent\textbf{本節のまとめと位置づけ．}
HFE を本章（§8）で定式化した理由は，CCD の Hermite データを基盤とする場延長理論が
PPE 離散化（§~\ref{sec:pressure}）および Balanced-Force 条件（§~\ref{sec:balanced_force}）と
密接に関連するためである．
§~\ref{sec:split_ppe_motivation} で議論する分相 PPE 解法では，界面上の圧力ジャンプ処理に
HFE が不可欠となり，高密度比シミュレーション（$\rho_l/\rho_g > 10$）への
推奨アップグレードパスとして位置づけられる．

% ─────────────────────────────────────────────────────────────────────────────
% backward-compat: rename to HFE in final draft
\phantomsection\label{eq:gfm_rhs_correction}%
\phantomsection\label{eq:gfm_ccd_div}%
% ─────────────────────────────────────────────────────────────────────────────

% ═══════════════════════════════════════════════════════════════════════════════
